% !PRINTED 419
\weberpage{419}
\begin{center}
\S.\,111.

\gesperrt{Die Resolventen der Ikosaëdergleichung.}
\end{center}

Die Formulirung des Problems, die wir am Schluss des vorigen Paragraphen gegeben haben, führt uns auf die Frage nach den verschiedenen Resolventen der Ikosaëdergleichung. Jedem Theiler der Ikosaëdergruppe entspricht eine solche Resolvente, deren Grad gleich dem Index des Theilers, also immer ein Theiler von 60 ist. Nun enthält die Ikosaëdergruppe (\S.\,59):

\begin{tabular}{@{}ll c l@{}}
1) Tetraëdergruppen (\S.\,56): & Resolventen & $5^{\mathrm{ten}}$ & Grades \\
2) Diëdergruppen $D_5$ (\S.\,55): & \textquotedbl & $6^{\mathrm{ten}}$ & \textquotedbl \\
3) Diëdergruppen $D_3$: & \textquotedbl & $10^{\mathrm{ten}}$ & \textquotedbl \\
4) Cyklische Gruppen $C_5$: & \textquotedbl & $12^{\mathrm{ten}}$ & \textquotedbl \\
5) Vierergruppen $D_2$: & \textquotedbl & $15^{\mathrm{ten}}$ & \textquotedbl \\
6) Cyklische Gruppen $C_3$: & \textquotedbl & $20^{\mathrm{sten}}$ & \textquotedbl \\
7) Cyklische Gruppen $C_2$: & \textquotedbl & $30^{\mathrm{sten}}$ & \textquotedbl \\
\end{tabular}

Da die Ikosaëdergruppe einfach ist, sind alle diese Gleichungen \gesperrt{Totalresolventen} von einander; wir beschränken uns hier auf die Betrachtung der wichtigsten unter ihnen.

Die Resolventen niedrigsten, nämlich $5^{\mathrm{ten}}$ Grades sind die zu der Tetraëdergruppe gehörigen.

Um sie zu finden, müssen wir eine Function von $\vartheta$ suchen, die durch die Substitutionen der Tetraëdergruppe ungeändert bleibt, während sie in der Ikosaëdergruppe fünfwerthig ist.

Wir nehmen die Ikosaëdergruppe in der in \S.\,58, (22) aufgestellten Normalform, und nehmen die darin enthaltene Tetraëdergruppe $Q$ (\S.\,59) heraus, die wir über der Vierergruppe
\weq{1}{1, \psi, \chi, \psi\chi}
construirt haben, worin, wenn $\varepsilon$ eine imaginäre fünfte Einheitswurzel ist, $\psi$ und $\chi$ die Bedeutung haben:
\weq{2}{\psi(x) = \dfrac{-1}{x}, \quad \chi(x) = \dfrac{\omega x+1}{x-\omega}, \quad \omega = \varepsilon + \varepsilon^{-1}.}

Setzen wir $\Theta(x) = \varepsilon x$, so wird die ganze Ikosaëdergruppe
\weq{3}{P = Q + Q\Theta + Q\Theta^2 + Q\Theta^3 + Q\Theta^4,}
und
\weq{4}{Q_h = \Theta^{-h} Q \Theta^h}
sind die zu $Q$ conjugirten Tetraëdergruppen.
